% include file for chapter 4

\newpage \section{VÝPOČETNÍ ČÁST} \label{comp}
Stěžejní částí této práce je počítačový program vytvořený v prostředí Matlab, který se skládá se z několika oddělených, avšak na sebe navazujících částí, kterými jsou: zpracování dat LLE na rovnovážné funkce, patrový model extraktoru a model přenosový.
V této kapitole budou popsány principy fungování programu a způsob výpočtů.

\subsection{Zpracování rovnovážných dat} \label{compLLE}
Tato část programu slouží ke zpracování naměřených dat LLE na funkce popisující rovnováhu ve tvaru $u = f(w, t)$.
Jak již bylo rozebráno v kapitole~\ref{LLEtheory}, vzhledem k malému koncentračnímu rozmezí je možné použít empirickou rovnici namísto modelů aktivitních koeficientů.
V daném rozsahu $w$ a $t$ plně postačuje funkce, která je polynomem třetího řádu vůči $w$ (avšak bez absolutního členu, neboť funkce musí procházet bodem $(0,0)$) a polynom prvního řádu (lineární funkce) vůči $t$, jak je popsáno v rovnici~\ref{eqn:eqFun0}. Tato funkce má 5 parametrů $p_1$--$p_5$.
\begin{align}
\label{eqn:eqFun0}
	u &= (p_1 w^3 + p_2 w^2 + p_3 w)(p_4 t + p_5) \\
\label{eqn:eqFun}
	u &= p_1 w^3 t + p_2 w^2 t + p_3 w t + p_4 w^3 + p_5 w^2 + p_6 w
\end{align}
V tomto tvaru jsou však parametry na sobě jasně závislé.
Numerický algoritmus na nelineární regresi, který je v prostředí Matlab poskytován, potom hledá optimum obtížně a nespolehlivě, proto byla funkce použita v jiném tvaru, který je popsán v rovnici~\ref{eqn:eqFun}. Tato funkce má 6 parametrů $p_1$--$p_6$.

Do regrese by sice bylo možné zahrnout obsah síranu amonného ve vodné fázi jako třetí nezávislou proměnnou, regrese by tím však byla méně spolehlivá.
Obsah SA lze přitom považovat za zaručený, z důvodů popsaných v kapitole~\ref{CHON} jej navíc ani není možné příliš měnit, proto byly jako nezávislé proměnné zahrnuty pouze $w$ a $t$.

Výše popsaná nelineární regrese je implementována ve funkci \textit{fitEqCurves.m}, kde je využita vestavěná funkce \textit{lsqcurvefit} nabízená prostředím Matlab, aby byla získána sada parametrů pro každou dvojici fází (vodná fáze s daným obsahem SA a vybrané rozpouštědlo).
Po úspěšném proložení regresních křivek lze spouštět funkci \textit{getEqFun.m}, která do rovnovážné funkce dosadí zadanou teplotu a získanou funkci $u = f(w)$ tabeluje jako vektory $u$ a $w$.
Tabelace funkce je výhodná pro výpočetně nenáročnou implementaci funkce inverzní, tedy $w = f^{-1}(u)$, která je hojně využívána v modelech extraktorů.

Pro vizualizaci naměřených dat byla vytvořena ještě funkce \textit{plotLLEdata.m}, která do grafu závislosti $u(w)$ vykreslí pro zadanou dvojici fází všechna data LLE, která jsou seskupena do sad dle teploty při měření.
Tyto grafy budou ukázány v kapitole~\ref{resLLE}.

\subsection{Implementace patrového modelu} \label{compStage}
Rovnovážné funkce využívá patrový model extraktoru, který pracuje s konceptem teoretického patra (TP).
Model ze zadané bilance určí počet teoretických pater (NTP) pomocí McCabe"-Thielovy metody~\cite{McCabe1925} modifikované pro kapalinovou extrakci, jak ji popisuje T.~W.~Evans~\cite{Evans1934}.
Kromě určení NTP je také vypočtena minimální hodnota hmotnostního toku solventu ($\dot{m}_\mathrm{S,min}$), což je klíčová informace pro navrhování procesu, neboť počet TP potřebných pro separaci se zvyšuje s klesajícím tokem solventu a při hodnotě $\dot{m}_\mathrm{S,min}$ limitně dosahuje nekonečna, při nižších hodnotách pak separace ani není možná.
Centrální funkcí tohoto modelu je \textit{extractor.m}, kde se nejprve dopočte hmotnostní bilance, načež jsou určeny hodnoty $\dot{m}_\mathrm{S,min}$ a NTP.

Bilanci je díky nemísitelnosti fází možné vyjádřit pomocí relativních hm. zlomků $W$ a hm. toků obou fází bez CPL, tedy toku samotného síranového roztoku $\dot{m}_\mathrm{SA}$ a toku samotného solventu $\dot{m}_\mathrm{S}$ (ten je shodný s tokem vstupujícího proudu solventu, neboť obsah CPL je v ní nulový).
Z těchto šesti bilančních veličin musí být právě jedna libovolná veličina ponechána nedefinovaná, aby mohla být vypočtena a zajištěno tak, že rovnice~\ref{eqn:massBalance} platí.
\begin{equation}
\label{eqn:massBalance}
	\dot{m}_\mathrm{SA} W_\mathrm{F} + \dot{m}_\mathrm{S} W_\mathrm{S} = \dot{m}_\mathrm{SA} W_\mathrm{R} + \dot{m}_\mathrm{S} W_\mathrm{E}
\end{equation}
Algoritmus na výpočet $\dot{m}_\mathrm{S,min}$ funguje obdobně jako dobře známá metoda půlení intervalu, nejprve jsou tedy zvoleny dvě hodnoty toku solventu -- velmi nízká hodnota $\dot{m}_1$ a velmi vysoká hodnota $\dot{m}_2$, přičemž se očekává, že hodnota $\dot{m}_\mathrm{S,min}$ bude ležet mezi těmito krajními hodnotami.
Poté je zahájen cyklus, kdy je při každé iteraci určena hodnota $\dot{m}_3$ jako aritmetický průměr hodnot $\dot{m}_1$ a $\dot{m}_2$ a následně je zjištěno, zda-li má pracovní přímka průsečík s rovnovážnou křivkou pro hodnoty $\dot{m}_1$, $\dot{m}_2$ a $\dot{m}_3$.
Pokud je průsečík nalezen pro $\dot{m}_1$ a $\dot{m}_3$, bude v další iteraci platit $\dot{m}_1 = \dot{m}_3$, zatímco pokud je průsečík nalezen pouze pro $\dot{m}_1$, bude platit $\dot{m}_2 = \dot{m}_3$.
Tyto iterace jsou prováděny tak dlouho, dokud interval $\dot{m}_1$--$\dot{m}_2$ není velice úzký ($\dot{m}_\mathrm{S} \cdot 10^{-4}$), protože s takovou přesností je hodnota $\dot{m}_\mathrm{S,min}$ určena.
Jako výslednou hodnotu lze použít $\dot{m}_3$.

Protože tato funkce poskytuje jen strojový výstup, byla vytvořena také funkce \textit{extractorPlot.m}, která přehledně vypíše důležité výsledky do konzole prostředí Matlab a vykreslí graf závislosti $u(w)$, kde je vynesena rovnovážná křivka, pracovní úsečka a úsečky znázorňující jednotlivá teoretická patra.
Pro usnadnění parametrické studie byla vytvořena ještě funkce \textit{extractorFeed.m}, která umožní vybranou vstupní veličinu měnit v definovaném rozmezí, výsledkem je pak tabulka závislosti NTP a $\dot{m}_\mathrm{S,min}$ na vybraném parametru.

\subsection{Implementace přenosového modelu} \label{compTransfer}
Jak již bylo dříve zmíněno, přenosový model při předpovědi chování extraktoru vychází ze znalostí základních vlastností látek a systému.
Relevantní vlastnosti látek jsou hustota, dynamická viskozita a difuzivita kaprolaktamu v dané látce, vlastnostmi systému jsou rovnovážná funkce $u = f(w)$ a mezifázové napětí ($\gamma$).
Vlastnosti látek byly exportovány z databáze simulačního programu Aspen Plus V10~\cite{Aspen} v širokém rozmezí teplot, aby bylo možné jejich hodnoty dle potřeby interpolovat za zvolené teploty.
Rovnovážná funkce tabelovaná jako vektor $u$ a $w$ je získána spuštěním funkce \textit{getEqFun.m} popsané v kapitole~\ref{compLLE}.
Naměřené hodnoty mezifázového napětí byly zpracovány na funkci $\gamma(w)$, k čemuž postačuje regrese naměřených dat pomocí polynomu druhého řádu (viz kapitolu~\ref{resgamma}).
Do této funkce je dosazena střední hodnota $w$ v profilu extraktoru, aby byla získána střední hodnota $\gamma$, se kterou je dále pracováno.

Přenosový model slouží k výpočtu koncentrací CPL v obou fázích v závislosti na výškové souřadnici extraktoru, a může být realizován buďto jako výpočet návrhový, kdy je vypočtena výška $L$ potřebná na dosažení požadované hm. bilance, nebo jako výpočet simulační, kdy je $L$ specifikováno a model předpovídá míru separace v extraktoru, čili doplní hm. bilanci.
Také průměr kolony $d_\mathrm{C}$ může být buďto specifikován (simulační výpočet) nebo vypočten z meze zahlcení (návrhový výpočet).
Libovolná kombinace těchto cest výpočtu je možná, a na tom pak závisí způsob běhu programu.

Celý tento model je implementován ve funkci \textit{design.m}, který nejprve zpracuje zadání a získá příslušná fyzikální data pro zadané fáze.
Při návrhu $L$ je bilance specifikována tak, jak bylo popsáno pro patrový model -- jedna libovolná bilanční veličina musí zůstat nedefinovaná a bilance je doplněna zavoláním funkce \textit{extractor.m}, která zároveň vypočte NTP a $\dot{m}_\mathrm{S,min}$.
V případě simulace $L$ musí být specifikovány všechny bilanční veličiny kromě výstupních koncentrací $w_\mathrm{R}$ a $w_\mathrm{E}$, protože míra separace je předmětem výpočtu, je tedy neznámá.
Program se dále větví na podprogramy, které pomocí empirických rovnic určí potřebné veličiny charakterizující přenos hmoty v extraktoru -- k tomu slouží funkce \textit{designPacked.m}, \textit{designRotDisc.m} a \textit{designKarr.m}, které budou podrobně vysvětleny v kapitole~\ref{compInsfun}.
Nejprve jsou zavolány pro výpočet meze zahlcení, která je buďto využita pro návrh $d_\mathrm{C}$, nebo, pokud byla tato hodnota specifikována, je mez zahlcení jen informativní.
Podruhé (při známých rozměrech kolony) jsou funkce zavolány pro určení všech ostatních potřebných veličin.
Následuje samotná implementace diferenciálního modelu, který bude podrobně popsán v kapitole~\ref{compDiff}.
Výsledkem modelu je při návrhu $L$ jeho určená hodnota, při specifikaci $L$ pak bilanční veličiny určené jak pro model s pístovým tokem, tak model s axiální disperzí.
V případě simulace $L$ jsou teprve nyní známy výstupní koncentrace, takže může být zavolána funkce \textit{extractor.m} na určení hodnot NTP a $\dot{m}_\mathrm{S,min}$.

\subsubsection{Řešení diferenciálních rovnic} \label{compDiff}
Jádrem výpočtu je matematický model, který popisuje extraktor jako soustavu dvou nelineárních diferenciálních rovnic druhého řádu, jak již bylo podrobně popsáno v kapitole~\ref{the_Math}.

S.~Hartland a A.~Kumar~\cite{HartlandKumar1997} uvádějí analytické řešení diferenciálních rovnic při návrhovém výpočtu -- řešením modelu s pístovým tokem je přímé vyjádření $L$, zatímco řešením modelu s axiální disperzí je rovnice transcendentní pro $L$, z níž lze $L$ určit numerickým algoritmem na hledáním kořene.
Analytická řešení ovšem platí pouze za předpokladu lineární rovnováhy, což, jak bude ukázáno v kapitole~\ref{resLLE}, v našem případě neplatí, rovnovážná funkce je totiž značně nelineární.
Autoři uvádí, že při nelineární rovnováze diferenciální rovnice nemají analytické řešení a je nutné je řešit numericky.

K numerickému řešení diferenciálních rovnic s okrajovou podmínkou Matlab nabízí funkci \textit{bvp5c}, ta pro svůj běh potřebuje zaprvé funkci popisující diferenciální model prvního řádu, která ze závislých a nezávislých proměnných vypočítá derivace závislých proměnných, zadruhé funkci definující počáteční podmínky jakožto rovnice v anulovaném tvaru, zatřetí pak funkci sloužící jako počáteční odhad řešení.
V případě návrhového výpočtu je třeba zahrnout i počáteční odhad výšky extraktoru.

Protože algoritmy řešení okrajových problémů nebývají zcela robustní a řešení může často divergovat, je třeba poskytnout dobrý počáteční odhad řešení.
Model s axiální disperzí se bez kvalitního odhadu ukázal být velmi nestabilní, proto je nejprve numericky vypočten model s pístovým tokem, jehož řešení je snazší a zcela přitom postačuje jednoduchý odhad lineárních profilů koncentrace.
V případě návrhového výpočtu, kdy jsou koncentrace vstupních i výstupních proudů známy, jsou lineární profily od $c_\mathrm{F}$ do $c_\mathrm{R}$, resp. od $c_\mathrm{E}$ do $c_\mathrm{S}$, v případě simulačního výpočtu zcela postačí dosadit namísto $c_\mathrm{R}$ nulu a namísto $c_\mathrm{E}$ 1~\%.
Pro odhad výšky, je-li třeba, postačuje použít výšku teoretického patra 0,5~m, která se vynásobí počtem TP získaným zavoláním patrového modelu.
Diferenciální rovnice prvního řádu i jejich počáteční podmínky byly již popsány v rovnicích~\ref{eqn:diff1}, \ref{eqn:bound1_1}, \ref{eqn:bound1_2} a~\ref{eqn:boundExtra}.

Nalezené řešení pro model s pístovým tokem má využití pro sestavení počátečního odhadu pro model s axiální disperzí, ten je však nutné nejprve upravit ze soustavy dvou dif. rovnic druhého řádu, jak jsou uvedeny v rovnici~\ref{eqn:diff2}, na soustavu čtyř rovnic prvního řádu.
To je možné provést zavedením derivací koncentrací jako dvou dalších závislých proměnných: $c'_x = \D c_x / \D z$ a $c'_y = \D c_y / \D z$.
Diferenciální model ve formě připravené pro numerický řešič je rozveden v soustavě rovnic~\ref{eqn:diff24}.

\begin{align}
\label{eqn:diff24}
	\frac{\D c_x}{\D z} &= c'_x \\
	\frac{\D c_y}{\D z} &= c'_y \nonumber \\
	\frac{\D c'_x}{\D z} &=  Pe_x c'_x + Pe_x T_x (c_x - c_y^*) \nonumber \\
	\frac{\D c'_y}{\D z} &= -Pe_y c'_y - Pe_y T_x (c_x - c_y^*) \frac{Q_x}{Q_y} \nonumber
\end{align}
Protože řešení modelu s pístovým tokem je ve formě vektorů hodnot $c_x$ a $c_y$ odpovídajících vektoru $z$, odhadovou funkci pro model s axiální disperzí lze realizovat jako algoritmus, který pro požadované $z$ nalezne dva sousedící body existujícího řešení, mezi nimiž požadované $z$ leží.
Hodnoty $c_x$ a $c_y$ jsou lineárně interpolovány, zatímco hodnoty $c'_x$ a $c'_y$ jsou vypočteny běžnou dvoubodovou formulí pro numerickou derivaci.
Při návrhovém výpočtu je součástí řešení modelu s pístovým tokem také vypočtená výška extraktoru -- její dvojnásobek je následně použit jako nový počáteční odhad výšky.

Hlavním výsledkem modelu s axiální disperzí je pak řešení diferenciálního modelu, tedy funkce $(c_x, c_y, c'_x, c'_y) = f(z)$ ve formě vektorů hodnot.
Při návrhovém výpočtu je klíčovým výsledkem $L$, při simulačním výpočtu zase hmotnostní bilance extraktoru.

\subsubsection{Kvantifikace parametrů přestupu hmoty} \label{compInsfun}
V kapitole~\ref{insfun} byly podrobně popsány vztahy, pomocí nichž lze kvantifikovat veličiny charakterizující přestup hmoty v extraktoru, kterými jsou specifická mezifázová plocha $a$, koeficient prostupu hmoty $K_x$ a koeficienty axiální disperze $E_y$ a $E_x$.
V případě návrhového výpočtu je také vypočtena plocha průřezu kolony $A$ (resp. průměr kolony $d_\mathrm{C}$) vyplývající z meze zahlcení $v_\mathrm{C,F}$.
Mez zahlcení je však vypočtena i v případě simulačního výpočtu, neboť se jedná o velmi důležitou provozní vlastnost.

Výpočty se zcela liší pro výplňovou kolonu, extraktory s rotujícími míchadly a Karrův extraktor, a to nejen tvarem rovnic, ale i posloupností výpočtů.
Ve všech případech jsou však vypočteny následující veličiny jakožto mezivýsledky: střední velikost kapek $d_{32}$, zádrž dispergované fáze $\phi$ a rychlost kluzu $v_\mathrm{s}$, která z $\phi$ vyplývá dle rovnice~\ref{eqn:vSlipCalc}.
Hodnota $a$ pak přímo vyplývá z hodnot $d_{32}$ a $\phi$ dle rovnice~\ref{eqn:aint}.

Některé dříve uváděné rovnice nejsou explicitní pro veličinu, k jejímuž výpočtu slouží.
V takových případech byla rovnice převedena na anulovaný tvar a řešena pomocí \textit{fzero}, což je vestavěná funkce, kterou Matlab nabízí k hledání kořene.

Řešení výplňové kolony (viz kapitolu~\ref{insfunPacked}), které sestavil J."-P.~Duroudier~\cite{Duroudier2016}, je poměrně přímočaré a mez zahlcení lze přitom vypočítat odděleně od zbytku výpočtu.
Před samotným výpočtem je provedena kontrola, zda-li zvolený rozměr výplně splňuje dříve popsaná omezení, aby bylo možné zaručit spolehlivost dalších výpočtů.
Pak je nejprve vypočtena $v_\mathrm{C,F}$, následuje $\phi$, $d_{32}$ a empirický výpočet $v_\mathrm{s}$ pomocí $\phi$.
K dispozici je však také přímý geometrický výpočet, jako výsledná hodnota rychlosti kluzu proto poslouží průměrná hodnota obou metod výpočtu.
Pomocí Reynoldsova čísla v kontinuální fázi $Re_\mathrm{C}$, Schmidtových čísel pro obě fáze ($Sc$) a předchozích mezivýsledků je vypočten koeficient přestupu pro disp. fázi $\beta_\mathrm{D}$, Sherwoodovo číslo pro kont. fázi $Sh_\mathrm{C}$ (a z něj vyplývající $\beta_\mathrm{C}$) a Pécletova čísla pro kont. i disp. fázi.
Pécletovo číslo pro kont. fázi je přitom omezeno maximální hodnotou závisející pouze na $Re_\mathrm{C}$, kterou údajně nemůže překročit, ta je proto použita v případě, že by ji standardně vypočtená hodnota převyšovala.
Tento postup je implementován ve funkci \textit{designPacked.m}, jako užitečná pomůcka byla vytvořena také funkce \textit{packing.m}, která program prováže s databází předdefinovaných výplní, kterou J."-P.~Duroudier uvádí v příloze.

Řešení pro RDC či Kühniho extraktor (viz kapitolu~\ref{insfunRotDisc}), které sestavili S.~Hartland a A.~Kumar~\cite{HartlandKumar1997}, mají společný postup řešení, byť někdy využívají odlišné tvary rovnic.
Při návrhovém výpočtu program začíná iterativním cyklem, poněvadž měrný disipovaný výkon $\psi$ velmi strmě závisí na rozměru kolony, a tedy i rozměru míchadla.
Počáteční odhad je pak dle doporučení autorů definován jako $v_\mathrm{C} = 4 \cdot 10^{-3}~\mathrm{m/s}$, zatímco při simulačním výpočtu jsou hodnoty $d_\mathrm{C}$ i $v_\mathrm{C}$ od počátku známy.
Následuje výpočet ostatních geometrických veličin, výpočet $\psi$ a z něj vyplývající $d_{32}$.
Pro určení $v_\mathrm{C,F}$ je nejprve vypočtena zádrž při zahlcení $\phi_\mathrm{F}$, která z $d_{32}$ vychází.
Při návrhovém výpočtu je z hodnoty $v_\mathrm{C,F}$ určena nová hodnota $v_\mathrm{C}$ a cyklus se opakuje až do konvergence, kdy relativní rozdíl mezi $d_\mathrm{C}$ činí méně než $10^{-4}$.
Po ukončení cyklu je vypočtena $\phi$, $Sh_\mathrm{C}$ pro určení $\beta_\mathrm{C}$, přímá korelace $\beta_\mathrm{D}$ a hodnoty $E$ pro obě fáze.
Tento postup je implementován ve funkci \textit{designRotDisc.m}.

Řešení Karrova extraktoru (viz kapitolu~\ref{insfunKarr}), které bylo v této práci sestaveno za použití různých literárních zdrojů~\cite{Baird1984,Kumar1999,Baird1976,Baird1979}, je opět iterativní, a to kvůli potřebě znát střední hustotu $\bar{\rho}$ pro výpočet $\psi$.
Nejprve je tedy odhadnuta zádrž $\phi = 0,05$ a v každé iteraci je vypočteno nejprve $\psi$ a z něj vyplývající $d_{32}$.
Rovnice~\ref{eqn:KarrVC_phi}, která popisuje vztah mezi rychlostí kont. fáze $v_\mathrm{C}$ a $\phi$, je poprvé využita na výpočet $v_\mathrm{C,F}$ dosazením vypočtené hodnoty $\phi_\mathrm{F}$ (zádrž při zahlcení) a podruhé na určení $\phi$ pomocí $v_\mathrm{C}$ (ta vyplývá buďto ze zadaných rozměrů kolony, nebo vypočtené $v_\mathrm{C,F}$).
Cyklus se poté opakuje, a vzhledem k tomu, jak malý je vliv $\phi$ na $\bar{\rho}$, bývá velice rychle dosažena konvergence (kdy relativní rozdíl $\phi$ činí méně než $10^{-4}$).
Po ukončení cyklu je pak proveden výpočet $Sh_\mathrm{C}$, $Sh_\mathrm{D}$ a přímá korelace $E_\mathrm{C}$.
Protože korelace $E_\mathrm{D}$ nebyla v literatuře nalezena, je jako aproximace předpokládána rovnost $Pe_\mathrm{C} = Pe_\mathrm{D}$.
Tento postup je implementován ve funkci \textit{designKarr.m}.
